\subsection{Text Overlap Generation (TinyStories)}
\label{sec:text_overlap}

\paragraph{Task Description.}
The text overlap experiment evaluates overlap-aware span bridging using
\texttt{roneneldan/TinyStories} (20 ingested samples, 40\% RIGHT holdout).
TinyStories was chosen for its vocabulary regularity: stories share canonical
verbs, settings, and character archetypes, creating a dense web of chord
attractor bridges across samples.  This controlled overlap is precisely
what makes the boundary detection hypothesis testable: the system should
identify structural boundaries where the prompt's chord fingerprint overlaps
with a learned corpus span, and transition into the subsequent span
naturally.

\paragraph{Results.}
Across $N = 20$ test samples the mean weighted score was 0.548.

\paragraph{Assessment.}
The substrate correctly identified and exploited chord-level overlap at
story span boundaries, producing continuations that bridge naturally
into the corpus.  The high score validates the hypothesis that TinyStories'
regular structure creates strong attractor bridges.

Figure~\ref{fig:text_overlap_map} shows the trial outcome map.

\begin{figure}[htbp]
  \centering
  \InputIfFileExists{text_overlap_map.tex}{}{}
  \caption{Text overlap trial map ($N = 20$ TinyStories samples,
    40\% RIGHT holdout).  Chord bridges across story span boundaries.}
  \label{fig:text_overlap_map}
\end{figure}
